% Copyright 2026 Alan Szmyt. All rights reserved pending publication license.
\section{Limitations}
\label{sec:limitations}

This manuscript establishes a design and a prospective feasibility protocol,
not treatment efficacy, clinical safety, successful affect regulation,
learnability of a personal mapping, universal audio--state relationships, or a
neurological mechanism. Its executable evidence is presently limited to a
deterministic synthetic session and publication checks at a named source
revision. The real-model, two-rate rendering, private-data, and formal-study
paths remain incomplete or blocked. A requirement stated in an architecture or
protocol is therefore classified below as \emph{specified}, not as an effective
mitigation. Table~\ref{tab:risk-mitigation-status} summarizes the current
controls beside their residual risks and blocking gates.

\subsection{Construct and measurement validity}
\label{sec:limitations-construct-validity}

The moment representation is a bounded account produced through a particular
interface; it is not the person's true psychological or physiological state.
Self-report can preserve experience that an external classifier cannot observe,
but wording, timing, scale anchors, recall, current arousal, and the wish to make
the session coherent may all shape a response. Compact affect scales provide a
measurement precedent \citep{betella2016slider}, yet that precedent does not
validate Antidote's desired-transition, mismatch, usefulness, harm, burden, or
aftereffect items. Likewise, optional physiology or behavior would be an
uncertain observation rather than an objective need: passive-monitoring
research reports substantial heterogeneity and missing-data concerns, and
electrodermal activity does not uniquely identify felt valence
\citep{deangel2022passive,damelio2025eda}.

Semantic intent, requested acoustic controls, measured realization, perceived
expression, felt response, immediate usefulness or harm, and later aftereffect
are kept in separate records. This reduces category errors; it does not prove
that any record validly or reliably measures its named construct. Perceived and
felt musical emotion may diverge \citep{zentner2008emotions}, and an analyzer's
tempo, key, loudness, or embedding estimate cannot decide what the audio meant
to the listener. Person correction and field-level missingness are consequently
necessary controls, but their usability, responsiveness, test--retest behavior,
and sensitivity to meaningful within-person change remain unvalidated.

The originating cathartic experience is one retrospective, researcher-authored
observation without a protocol, comparator, standardized measure, or independent
interpretation. It motivates the person-and-moment question but is not a Result,
an effect estimate, or evidence that emotional intensity was helpful. The
framework's deliberate separation of intensity, welcome or unwanted response,
harm, and later meaning must survive instrument review before any qualifying
human observation can support even a feasibility claim.

\subsection{Internal and statistical validity}
\label{sec:limitations-internal-validity}

The prospective H1 series is an unblinded self-study in which researcher,
developer, participant, analyst, and author roles overlap. That overlap creates
expectancy, demand, selective-attention, analytic-flexibility, and role-conflict
risks that provenance cannot remove. The generic, state-aware, and negotiated
personal conditions are also visibly different bundles: personal context,
semantic content, plan visibility, editing, and approval change together.
Neither a pairwise contrast nor exact randomization can isolate which component
produced a difference, and no contrast identifies a psychological or
neurological mechanism.

Baseline fluctuation, time of day, setting, prior activity, familiarity,
regression toward a person's usual state, repeated instrument use, novelty,
order, carryover, and learning may all covary with condition. The frozen
balanced schedule and within-block seed policy constrain some researcher
degrees of freedom, while near-event assessment can reduce retrospective delay
\citep{shiffman2008ema}. They do not create blinding, washout, exchangeability,
or a stable underlying affect process. A stopped or missing response may also be
informative rather than random, particularly when discontinuation follows
distress, burden, technical failure, or mismatch.

The target of 18 completed exposures under an absolute ceiling of 24 attempts
is appropriate only to the declared record-feasibility questions. The series
cannot estimate a population effect or support a clinical, safety, superiority,
or causal claim. SPENT and CENT provide useful protocol and complete-reporting
precedents for N-of-1 work \citep{porcino2020spent,vohra2015cent}; checklist
alignment does not repair weak measurement, multiplicity, bundled conditions,
carryover, or investigator conflict. Preserving every assignment, failure,
exposure, null or opposite-direction observation, adverse response, denominator,
missingness reason, and prespecified analysis is a transparency control. It is
not a guarantee of conclusion validity.

\subsection{External and ecological validity}
\label{sec:limitations-external-validity}

Any later H1 observation would describe one person under one frozen protocol,
model, adapter, device path, listening environment, semantic vocabulary, and
period of life. It would not establish transfer across people, moments,
languages, cultures, musical histories, sensory tolerances, hearing profiles,
devices, generators, or clinical contexts. Even within one person, a mapping
may change after familiarity, major events, medication or substance changes,
sleep disruption, illness, or a shift in what the same semantic term means.
Versioning can locate such change; it cannot make the relationship stationary.

The proposed local, self-guided setting is not interchangeable with a music
therapy relationship, laboratory procedure, guided psychedelic setting, or
ordinary music listening. The presence of headphones, an eye mask, a check-in,
or a structured journey may itself alter attention and expectancy. Reviews of
music-based digital therapeutics and music-neurofeedback systems report
heterogeneous methods and limited grounds for real-world generalization
\citep{venkatesan2026mdt,sayal2025musicloop}. Those limitations neither prove
that Antidote cannot work nor support transporting a favorable demonstration.
A later multi-participant protocol would need a new sampling rationale,
accessible procedure, device and setting qualification, and explicit tests of
which constructs and controls transfer.

\subsection{System, model, and continuity failures}
\label{sec:limitations-system-failures}

Every transformation in the proposed chain can fail while producing a
plausible-looking artifact. A context extractor may omit the relevant detail or
privilege the wrong one; a state estimate may become confidently wrong; a
semantic mixin may conflict with an exclusion; an adapter may ignore or
downgrade a control; an analyzer may mischaracterize the output; and a
generative model may reproduce, imitate, or introduce unwanted material. Human
correction constrains authority but does not ensure that a person can identify
every omission, understand every control, or review effectively while
distressed.

The two-rate architecture is also prospective. The checked-in application can
supervise a deterministic mock worker, verify a synthetic WAV, start deliberate
playback, stop, recover, and record a response. It does not yet implement or
qualify a real music model, uncertain state estimator, receding-horizon
controller, verified future-audio buffer, or deterministic continuity renderer.
Real generation can miss deadlines; supported controls may differ by adapter;
crossfades can hide a waveform discontinuity while leaving harmony, rhythm,
timbre, meaning, or pacing incoherent. Normative scheduling and loudness
references help define tests \citep{w3c2024webaudio,ebu2023r128}; they do not
establish seamlessness, perceived continuity, or an appropriate experience.

Model, code, analyzer, driver, and hardware drift may invalidate a prior result
even when a prompt appears unchanged. Immutable specifications, capability
negotiation, hashes, event replay, versioned analyzers, cancellation, fallbacks,
and complete failure accounting can make some divergence detectable. Provenance
records what the governed system says occurred; W3C PROV and workflow-run
packaging do not prove semantic correctness, authenticity, consent, ownership,
privacy, or benefit \citep{moreau2013prov,leo2024rocrate}. T1 must therefore
qualify control adherence, latency, buffer margin, continuity, integrity,
interruption, and recovery on the exact real-model and hardware manifest before
that path can inform H1. Fallbacks that dominate playback, unverifiable outputs,
missed deadlines, or unreliable stopping count against feasibility rather than
as invisible implementation details.

\AntidoteTable{risk-mitigation-status}

\section{Ethics statement}
\label{sec:ethics-statement}

No formal participant study, recruitment, human-subject data collection, ethics
approval, exemption, or independent determination is reported. Protocol
\texttt{ANT-PROT-FEAS-001} remains
\texttt{blocked-no-collection-authority}. Publication of this design and
protocol cannot activate it. Antidote is a research instrument, not a treatment,
diagnostic, medical device, crisis service, or substitute for professional care.

\subsection{Emotional safety and participant authority}
\label{sec:ethics-emotional-safety}

Music can be unwanted or harmful as well as meaningful. A theoretical model of
music-induced harm identifies affective, behavioral, cognitive, identity,
interpersonal, physical, and spiritual dimensions and treats harm as an
interaction among context, music, recipient, and delivery
\citep{silverman2020complicated}. Antidote must therefore never optimize tears,
intensity, duration, completion, engagement, or physiological change as a proxy
for benefit. Crying, calm, activation, numbness, discomfort, catharsis, and
distress require the person's interpretation in that moment and again, when
appropriate, at the declared later interval.

The design preserves explicit playback, skip, stop, correction, grounding, and
withdrawal authority and gives mismatch, unwanted intensity, harm, hearing
discomfort, serious or unresolved events, and later aftereffect their own
records. A stop must halt playback independently of model completion; a safety
halt must block automatic continuation and personal-model updating. These are
necessary requirements, not evidence that distress is prevented or resolved.
The prototype provides no clinical monitoring, emergency response, or reliable
way to distinguish tolerable challenge from escalating harm.

Before any human exposure, an independently reviewed safety plan must define
eligibility and exclusions, pre-session information, low-burden stopping,
grounding and follow-up, adverse-event classification, escalation contacts,
unresolved-event review, and discontinuation rules. High-risk crisis use,
unsupervised treatment claims, and guidance for medication or psychoactive-drug
sessions remain outside v0 authority. Refusal, stopping, non-response, or an
unfavorable account must not reduce access, be reframed as noncompliance, or be
used as a negative reward signal.

\subsection{Consent, privacy, and data governance}
\label{sec:ethics-consent-privacy}

Consent is action- and purpose-specific. Permission to inspect one source does
not authorize summarization, generation, playback, retention, learning,
synchronization, research analysis, export, or publication. The current
synthetic desktop path requires explicit confirmation for a manually entered,
session-scoped context projection and keeps the worker outside unrestricted
history. Historical journal or therapy-chat integration, revocation UI,
retention enforcement, and privacy-ready payload handling are not implemented.
An inspectable grant or projection also does not demonstrate comprehension,
voluntariness, or freedom from demand---a particular concern when researcher
and participant are the same person.

Local-first design can reduce required server transfer and keep an authoritative
record available on the person's device \citep{kleppmann2019localfirst}; it is
not perfect security. The current store is not ready for private participant
data. At-rest encryption, key creation and recovery, operating-system account
separation, backup policy, retention enforcement, auditable deletion,
tombstones, derived-projection invalidation, crash-log review, and a production
worker sandbox remain unresolved. Device compromise, malware, cloud-synced
folders, backups, screen capture, exports, and another local user can still
expose sensitive context or response data.

Public provenance presents a second tension: detailed lineage aids audit but can
increase linkability and re-identification risk. The reporting contract permits
only privacy-reviewed aggregate public envelopes and keeps row-level human
records, free text, private paths, stable identifiers, record hashes, and the
private evidence index outside the repository. Even a reconciled aggregate
commitment does not establish anonymity, consent, authenticity, truth, or
completeness. A person must be able to decline retention and export separately;
the exact deletion effect on events, payloads, projections, analyses, backups,
and already shared artifacts must be specified before collection. No public
release may proceed solely because a schema validates or a hash matches.

\subsection{Licensing, bias, accessibility, and hearing safety}
\label{sec:ethics-licensing-accessibility}

No real generator has been selected or qualified for the study. Before a model
or audio artifact can enter T1, H1, a release, or the magazine, review must cover
the exact code and weight licenses, redistribution and commercial restrictions,
training-data documentation and consent, required attribution, model-card
limitations, remote-code execution, recognizable reproduction or imitation,
third-party samples, output rights, and the intended distribution context.
Generative-audio research has historically discussed positive implications far
more often than negative ones and identifies copyright, agency, bias, privacy,
misuse, labor, access, and resource concerns as material review categories
\citep{barnett2023ethical}. That review does not prove that a particular model
or output violates rights; conversely, an open weight license cannot clear its
training data or every generated artifact.

Training corpora, control vocabularies, affect labels, and musical evaluation
may privilege well-represented genres, languages, cultures, instruments, and
listener assumptions. A person-extensible semantic layer may reduce pressure to
adopt one vocabulary, but adapter mappings and model outputs can still flatten
or appropriate meaning. The prototype has automated keyboard and semantic
checks; it has not established usability with screen readers, reduced-motion or
low-vision settings, hearing differences, motor constraints, cognitive load,
trauma-related sensitivities, or users outside the developer context. More
knobs may increase agency for one person and burden or suggestion effects for
another. Accessibility and cultural review must precede recruitment rather than
be inferred from a nominally local interface.

File-level integrated loudness and true-peak constraints can bound some
digital-output errors but do not determine sound pressure at the ear. The
WHO--ITU safe-listening standard treats exposure as level over duration and
recommends dosimetry, accessible information, and volume-limiting options
\citep{whoitu2019safelistening}. Antidote currently has no calibrated device
path, at-ear measurement, cumulative-dose accounting, or proof that its lowest
volume is appropriate. Human exposure therefore requires device-aware
qualification, conservative defaults, continuously available volume and stop
control, warnings that do not rely on sound alone, and a response path for
tinnitus, pain, dizziness, or hearing discomfort. None of those controls alone
establish hearing safety.

\subsection{Independent review and collection gate}
\label{sec:ethics-review-gate}

The project has not recorded an institutional review board, research ethics
committee, or other independent determination, and this manuscript does not
decide whether a future self-study is exempt or outside a particular regime.
That determination depends on the final protocol, institution, jurisdiction,
recruitment, data, claims, and intended use. Self-use, local execution, public
source code, or the absence of payment does not by itself resolve research
ethics. WHO health-AI guidance provides a relevant normative foundation for
autonomy, well-being, transparency, accountability, inclusiveness, and
responsiveness \citep{who2021ethics}; alignment with those principles is not an
approval or safety certification.

Before H1 can change state, independent review must document the applicable
ethics determination and approve or otherwise disposition the exact protocol,
role conflict, consent process, eligibility and exclusion rules, emotional and
hearing-safety plan, adverse-event and crisis escalation, compensation if any,
privacy and data-management plan, retention and deletion behavior, public-export
boundary, accessibility procedure, device qualification, model and output
rights audit, analysis plan, and stopping rules. The separate real-model/T1,
privacy, consent, independent-review, assignment, instrument, analysis, and
pilot stop/go gates must each be satisfied by named evidence; no gate is
inherited from the paper or another project. Any material amendment requires a
new version and renewed review before collection continues.

Until those records exist, formal collection, recruitment, therapeutic use,
autonomous personalization, sensor inference, and public release of private or
row-level response data remain prohibited. A later approval would be limited to
its exact scope and would not prove benefit, privacy, accessibility, or safety.
